% from ega0-1-fr.tex
\oldpage[0\textsubscript{I}]{67}

\begin{corollary}[7.3.3]
\label{0.7.3.3-fr}
Sous les hypothèses de (7.3.2), l'application canonique
$M\otimes_A\widehat A\to\widehat M$ est bijective, et le foncteur
$M\otimes_A\widehat A$ est exact en $M$ dans la catégorie des $A$-modules de
type fini~; par suite, le séparé complété $\mathfrak J$-adique $\widehat A$
est un $A$-module plat (6.1.1).
\end{corollary}

Notons d'abord que dans la catégorie des $A$-modules de type fini,
$\widehat M$ est un foncteur \emph{exact} en $M$. Soit en effet
\[
  0\to M'\xrightarrow{u}M\xrightarrow{v}M''\to0
\]
une suite exacte~; on sait déjà que
$\widehat v:\widehat M\to\widehat{M''}$ est surjectif (7.3.1)~; d'autre part,
si $i$ est l'homomorphisme canonique $M\to\widehat M$, il résulte du th. de
Krull que l'adhérence dans $\widehat M$ de $i(u(M'))$ s'identifie au séparé
complété de $M'$ pour la topologie $\mathfrak J$-préadique~; donc
$\widehat u$ est injectif, et en vertu de (7.3.1) l'image de $\widehat u$ est
égale au noyau de $\widehat v$.

Cela étant, l'application canonique $M\otimes_A\widehat A\to\widehat M$
s'obtient en passant à la limite projective sur les applications
\[
  M\otimes_A\widehat A\to
  M\otimes_A(A/\mathfrak J^n)=M/\mathfrak J^nM.
\]
Il est clair que cette application est bijective lorsque $M$ est de la forme
$A^p$. Si $M$ est un $A$-module de type fini, on a une suite exacte
$A^p\to A^q\to M\to0$, d'où, en vertu de l'exactitude \emph{à droite} des
foncteurs $M\otimes_A\widehat A$ et $\widehat M$ (en $M$) dans la catégorie
des $A$-modules de type fini, le diagramme commutatif
\[
\begin{tikzcd}[column sep=large,row sep=large]
  A^p\otimes_A\widehat A \arrow[r] \arrow[d] &
  A^q\otimes_A\widehat A \arrow[r] \arrow[d] &
  M\otimes_A\widehat A \arrow[r] \arrow[d] &
  0 \\
  \widehat A^{\,p} \arrow[r] &
  \widehat A^{\,q} \arrow[r] &
  \widehat M \arrow[r] &
  0
\end{tikzcd}
\]
où les deux lignes sont exactes et les deux premières flèches verticales des
isomorphismes~; on en tire aussitôt la conclusion.

\begin{corollary}[7.3.4]
\label{0.7.3.4-fr}
Soient $A$ un anneau noethérien, $\mathfrak J$ un idéal de $A$, $M$, $N$ deux
$A$-modules de type fini~; on a des isomorphismes canoniques fonctoriels
\[
  (M\otimes_A N)^\wedge
  \xrightarrow{\sim}
  \widehat M\otimes_{\widehat A}\widehat N,
  \qquad
  \bigl(\operatorname{Hom}_A(M,N)\bigr)^\wedge
  \xrightarrow{\sim}
  \operatorname{Hom}_{\widehat A}(\widehat M,\widehat N).
\]
\end{corollary}

Cela résulte de (7.3.3), (6.2.1) et (6.2.2).

\begin{corollary}[7.3.5]
\label{0.7.3.5-fr}
Soient $A$ un anneau noethérien, $\mathfrak J$ un idéal de $A$. Les
conditions suivantes sont équivalentes~:
\begin{enumerate}
  \item[$a)$] $\mathfrak J$ est contenu dans le radical de $A$.
  \item[$b)$] $\widehat A$ est un $A$-module fidèlement plat (6.4.1).
  \item[$c)$] Tout $A$-module de type fini est séparé pour la topologie
  $\mathfrak J$-préadique.
  \item[$d)$] Tout sous-module d'un $A$-module de type fini est fermé pour la
  topologie $\mathfrak J$-préadique.
\end{enumerate}
\end{corollary}

Comme $\widehat A$ est un $A$-module plat, les conditions $b)$ et $c)$ sont
équivalentes, car $b)$ équivaut à dire que si $M$ est un $A$-module de type
fini, l'application canonique
$M\to\widehat M=M\otimes_A\widehat A$ est injective (6.6.1, $c)$). Il est
immédiat que $c)$ entraîne $d)$, car si $N$ est un sous-module d'un $A$-module
$M$ de type fini, $M/N$ est séparé pour la topologie $\mathfrak J$-préadique,
donc $N$ est fermé dans $M$. Montrons que $d)$ implique $a)$~: si
$\mathfrak m$ est un idéal maximal de $A$, $\mathfrak m$ est fermé dans $A$
pour la topologie $\mathfrak J$-préadique, donc
\[
  \mathfrak m=\bigcap_{p\geq0}(\mathfrak m+\mathfrak J^p),
\]
et comme $\mathfrak m+\mathfrak J^p$ est nécessairement égal à $A$ ou à
$\mathfrak m$, on a $\mathfrak m+\mathfrak J^p=\mathfrak m$ pour $p$ assez

